\section*{Bab \ref{ch:g2:caring-for-nature} --- Merawat Alam}

\begin{solution}{exo:g2:caring-for-nature:1}
Sekeliling yang menjadi sandaran sebuah \omterm{def:g1:living-or-not:living}{makhluk hidup}. Misalnya: udara
yang kita hirup, air yang kita minum, makanan yang kita makan ---
semuanya dari \omterm{def:g2:caring-for-nature:environment}{alam}.
\end{solution}

\begin{solution}{exo:g2:caring-for-nature:2}
Misalnya: kumbang di bawah kulit kayunya, kutu kayu di bawah batang
kayunya, lumut di atasnya, landak yang melewati musim dingin di
rongganya.
\end{solution}

\begin{solution}{exo:g2:caring-for-nature:3}
Misalnya: bawa pulang semua sampah; amati alih-alih mengumpulkan, dan
jangan pernah memetik seluruh rumpun; kembalikan setiap batu atau
batang kayu yang dibalikkan; tetaplah di jalan setapak di dekat tempat
yang rapuh; amati \omterm{def:g1:animals-around-us:animal}{hewan} dengan tenang dari jauh dan biarkan anaknya
tidak terganggu.
\end{solution}

\begin{solution}{exo:g2:caring-for-nature:4}
Karena \omterm{def:g2:caring-for-nature:environment}{alam} tidak punya truk sampah: apa yang kita tinggalkan tetap
tinggal --- sebuah tutup plastik selama bertahun-tahun yang amat
panjang, cukup panjang untuk menjerat atau meracuni \omterm{def:g1:animals-around-us:animal}{hewan} dan tetap
tergeletak di situ ketika anak-anak hari ini sudah dewasa.
\end{solution}

\begin{solution}{exo:g2:caring-for-nature:5}
Kalau dipetik sampai habis, rumpunnya tidak dapat membuat \omterm{def:g2:from-seed-to-plant:seed}{biji} ---
sehingga tahun depan tidak ada \omterm{def:g1:plants-around-us:parts}{bunga} aster di situ --- dan lebahnya
kehilangan \omterm{def:g1:plants-around-us:parts}{bunga} tempatnya mencari makan.
\end{solution}

\begin{solution}{exo:g2:caring-for-nature:6}
Mereka melapuk perlahan menjadi tanah gelap yang gembur. Cacing, kutu
kayu dan \omterm{def:g1:living-or-not:living}{makhluk hidup} yang terlalu kecil untuk dilihat mengerjakannya;
pada akhirnya kita memperoleh kompos --- tanah baik yang memberi makan
\omterm{def:g1:plants-around-us:plant}{tumbuhan} tahun depan.
\end{solution}

\begin{solution}{exo:g2:caring-for-nature:7}
Misalnya: sepiring kecil air pada gelombang panas, sudut kebun yang
liar dan tak dipangkas, \omterm{def:g2:from-seed-to-plant:seed}{biji} untuk burung pada musim dingin, menggali
sebuah kolam kecil.
\end{solution}

\begin{solution}{exo:g2:caring-for-nature:8}
Jawaban terbuka. Bungkus dan botol plastik biasanya menang. Tiap
jenisnya seharusnya pulang bersama pemiliknya --- lalu masuk ke tempat
pemilahan yang tepat.
\end{solution}

\begin{solution}{exo:g2:caring-for-nature:9}
Biarkan \omterm{ex:g2:animals-grow-up:frog}{berudunya} di kolam: kolamnya menyimpan segala yang mereka
butuhkan --- air kolam, makanan, ruang untuk menumbuhkan kaki --- dan
toples tidak. Kalau Emma ingin menolong, ia dapat mengamati mereka di
situ sepanjang musim semi, menjaga tepiannya tidak terinjak, dan
barangkali menolong membuat \omterm{def:g2:where-animals-live:habitat}{habitat} kolam baru seperti kolam sekolahnya.
\end{solution}

\begin{solution}{exo:g2:caring-for-nature:10}
Sejuta ``tidak berarti apa-apa'' ke arah yang buruk adalah sebuah
padang rumput yang terkubur tutup botol --- kerusakan kecil berlipat
ganda. Ke arah yang baik, sejuta perhatian kecil --- sampah yang
dibawa pulang, keran yang dimatikan, sampah yang dipilah ---
menjumlah dengan sama pastinya: \omterm{def:g2:caring-for-nature:environment}{alam} sebagian besar dirawat melalui
pilihan yang persis sekecil itu.
\end{solution}
